\newpage

% ######################################################################
% SECTION 24: REAL RUN POSTMORTEMS
% ######################################################################
\section{Real Run Postmortems --- Jo Asal Mein Hua}

\subsection{Yeh Section Kaisa Hai}

\begin{storybox}
Theory kafi ho gayi. Ab baat karte hain \textbf{asli runs} ki --- jo maine is project par execute kiye, asli API keys ke saath. Ye numbers aur observations \textbf{real} hain. Interview mein jab aap bologe ``maine pipeline chalayi, ye aaya, ye problem mili, aise fix kiya'' --- wo \textbf{experience} ka proof hai jo koi textbook nahi deta.
\end{storybox}

\subsection{Run 1: Solid State Battery Market --- End-to-End Success}

\begin{codestorybox}
\textbf{Run details (real):}
\begin{itemize}
\item \textbf{Topic:} ``solid state battery market outlook 2026''
\item \textbf{Duration:} \textasciitilde176 seconds (end-to-end, real API calls)
\item \textbf{Sources:} 15 web sources (3 sub-queries $\times$ 5 results) via Tavily
\item \textbf{Fact-check:} \textbf{PASSED on first pass} --- revision\_count = 0
\item \textbf{Output:} Publication-grade Markdown report
\item \textbf{Persistence:} ResearchSession \#1 in Django DB + ChromaDB memory
\end{itemize}

\textbf{Report structure (jo LLM ne banaya):}
\begin{enumerate}[label=\arabic*.]
\item Table of Contents
\item Executive Summary
\item Evidence-backed Key Findings table
\item Technical breakdown (chemistry table + CAGR table)
\item Consensus \& Contradictions section
\item Citations/References
\end{enumerate}
\end{codestorybox}

\begin{conceptbox}
\textbf{Is run se 3 lessons:}
\begin{enumerate}[label=\arabic*.]
\item \textbf{First-pass fact-check success} --- Analyst ka structured output (themes/contradictions/gaps) researcher ke raw data se itna consistent tha ki fact-checker ko koi critical gap nahi mili. Ye prompt design ka payoff hai.
\item \textbf{176s ka time budget} --- Tavily search (3 calls) + 4 LLM calls + embedding store. Sabse bada chunk: Tavily web calls aur ChromaDB embedding computation.
\item \textbf{Determinism nahi hai} --- LLM outputs vary karte hain; isliye benchmark multiple runs ka average mangta hai.
\end{enumerate}
\end{conceptbox}

\subsection{Run 2: RAG Benchmark --- Jab Multi-Agent Haar Gaya}

\begin{codestorybox}
\textbf{Run details (real):}
\begin{itemize}
\item \textbf{Topic:} ``Best practices for RAG in 2026''
\item \textbf{Duration:} \textasciitilde74.4s (baseline + pipeline + judge)
\item \textbf{Baseline:} Single LLM call (writer prompt, no research)
\item \textbf{Pipeline:} Full multi-agent flow (researcher $\rightarrow$ analyst $\rightarrow$ fact\_checker $\rightarrow$ writer)
\item \textbf{Judge:} LLM-as-judge, depth + verifiability (0-10)
\end{itemize}

\textbf{Result:}
\begin{longtable}{p{5cm} p{3.5cm} p{3.5cm}}
\toprule
\textbf{Metric (0-10)} & \textbf{Single-agent} & \textbf{Multi-agent} \\
\midrule
Depth & 9 & 8 \\
Verifiability & 8 & 7 \\
\bottomrule
\end{longtable}

\textbf{Verdict:} \textbf{SINGLE\_AGENT\_SUPERIOR} --- single prompt ne benchmark jeeta!
\end{codestorybox}

\begin{warningbox}
\textbf{Ye failure gold hai.} Multi-agent pipeline haar gayi kyunki us run mein Writer ko Analyst ka raw (excerpt-heavy) output mila aur report source digest jaisi ban gayi --- synthesized prose nahi. Interview mein ise \textbf{honestly} batana: ``Maine benchmark chalaya, ek run mein multi-agent haar gaya, ye cause tha, ye fix propose kiya'' --- ye dikhata hai ki aap apne system ki \textbf{weaknesses} jaante ho aur measurement seriously lete ho.
\end{warningbox}

\subsection{Why Did Multi-Agent Lose? --- Root Cause Analysis}

\begin{conceptbox}
\textbf{Handoff quality hypothesis:}

Multi-agent pipeline ki quality \textbf{chain ki weakest link} jitni hai:
\begin{enumerate}[label=\arabic*.]
\item Researcher 15 raw sources deta hai (URL + content snippets).
\item Analyst unhe \texttt{research\_data[:10]} tak limit karke \textbf{brief} banata hai --- par is run mein brief excerpt-heavy tha, synthesis kam.
\item Writer ko brief + top-8 sources milte hain. Agar brief already excerpt-heavy hai, to Writer ke paas \textbf{synthesize karne ki material} nahi hoti --- wo source list copy kar deta hai.
\end{enumerate}

\textbf{Contrast:} Solid-state run mein Analyst ne clean themes/contradictions structure banaya, isliye Writer ne behtareen report likhi. RAG run mein handoff weak tha.

\textbf{Conclusion:} Pipeline deterministic nahi hai; ek run ka verdict \textbf{statistically meaningless} hai. 3-5 topics ka average chahiye.
\end{conceptbox}

\subsection{Fix Proposals --- Agar Mein Aage Kaam Karun}

\begin{codestorybox}
\textbf{3 concrete fixes (implementation-ready):}

\begin{enumerate}[label=\arabic*.]
\item \textbf{Analyst few-shot:} Analyst prompt mein 1 example of ``theme with synthesis'' --- excerpt copy nahi, balki 2-line synthesized insight. Isse handoff quality stabilize hoti.
\item \textbf{Writer instructions:} ``Har section mein at least 3 synthesized sentences likho jo multiple sources ko combine karein; direct source excerpt mat copy karo.'' Negative instruction = prevention.
\item \textbf{Multi-topic benchmark:} Benchmark runner ko 5 fixed topics par chalao, average scores report karo. Single run ka verdict basura hai.
\end{enumerate}
\end{codestorybox}

\begin{interviewbox}
\textbf{Q: ``Agar aapko 1 month mile is project ko improve karne, top 3 cheezein kya karoge?''}

\textbf{Answer:} (1) \textbf{Evaluation infrastructure} --- 20-topic benchmark suite + regression tracking, kyunki bina measurement ke improvement nahi hota. (2) \textbf{Parallel researcher agents} --- 3 sub-queries ko alag parallel nodes mein chalao, latency 176s se 60s tak. (3) \textbf{Production hardening} --- rate limiting, async task queue (Celery), observability (LangSmith tracing), aur vector store versioning. Ye 3 system ko demo se product tak le jaate hain.
\end{interviewbox}

% ######################################################################
% SECTION 25: SYSTEM DESIGN
% ######################################################################
\section{System Design --- Poora Architecture Ek Nazar Mein}

\subsection{High-Level Architecture}

\begin{designbox}
\textbf{4-layer architecture:}

\begin{itemize}
\item \textbf{Presentation layer:} Streamlit (human UI) + MCP server (AI clients) + Django REST API (programmatic).
\item \textbf{Orchestration layer:} LangGraph StateGraph --- researcher/analyst/fact\_checker/writer nodes + conditional loop.
\item \textbf{Knowledge layer:} ChromaDB (research\_docs + past\_research) + Tavily web search.
\item \textbf{Model layer:} Groq API (GPT-OSS-120B) + sentence-transformers embeddings.
\end{itemize}

\textbf{Key design principle:} Har layer apni boundary maintain karta hai. Interface files (\texttt{app.py}, \texttt{mcp\_server.py}, \texttt{backend/api/views.py}) sirf \texttt{research\_graph} ko call karte hain --- pipeline implementation se decoupled.
\end{designbox}

\subsection{Data Flow --- Request Se Report Tak}

\begin{codestorybox}
\textbf{Full request journey (Streamlit example):}

\begin{enumerate}[label=\arabic*.]
\item User \texttt{topic} type karta hai, \textbf{Run} dabata hai.
\item \texttt{app.py} state banata hai: \texttt{\{'topic': topic, 'messages': []\}}.
\item \texttt{research\_graph.stream(state)} shuru --- har node ka event UI mein dikhta hai.
\item \textbf{Researcher:} 3 sub-queries (LLM) $\rightarrow$ 3 Tavily searches $\rightarrow$ \textasciitilde15 sources + RAG retrieval (5 docs + 3 past) $\rightarrow$ state update.
\item \textbf{Analyst:} top-10 sources + top-5 rag chunks $\rightarrow$ themes/contradictions/gaps brief (LLM).
\item \textbf{Fact\_checker:} brief + top-8 raw sources $\rightarrow$ VERDICT (LLM). Fail $\rightarrow$ researcher loop (max 3).
\item \textbf{Writer:} brief + top-8 sources $\rightarrow$ Markdown report (LLM) $\rightarrow$ ChromaDB memory store.
\item UI report render karta hai; \textbf{Download} buttons markdown/text.
\end{enumerate}
\end{codestorybox}

\subsection{Module Dependency Graph}

\begin{conceptbox}
\textbf{Dependency direction (acha design):}

\begin{lstlisting}[style=envstyle]
app.py / mcp_server.py / backend/api/views.py
        |            |                |
        v            v                v
      graph/workflow.py  <-----  (compiled singleton: research_graph)
        |
        v
   agents/*.py  ---------->  state/schema.py
        |
        +------------------>  config/setting.py
        |
        +------------------>  rag/vector_store.py
\end{lstlisting}

\textbf{Observations:}
\begin{itemize}
\item \texttt{state/schema.py} par \textbf{koi dependency nahi} --- pure data contract.
\item \texttt{config/setting.py} leaf node --- environment/config sab isse aata hai.
\item \texttt{rag/vector\_store.py} sirf config par depend karta hai --- reusable module.
\item Har agent sirf state + config + rag use karta hai --- \textbf{dependency inversion} (interface through state).
\end{itemize}
\end{conceptbox}

\subsection{Failure Modes and Resilience}

\begin{longtable}{p{4.2cm} p{4.2cm} p{5.2cm}}
\toprule
\textbf{Failure} & \textbf{Kya hota} & \textbf{Is project mein handling} \\
\midrule
Groq API down & LLM calls fail & try/except, fallback text with error message \\
Tavily key missing & No web search & Notice message in research\_data \\
ChromaDB missing & No RAG context & try/except pass (retrieval empty) \\
LLM returns garbage & Bad report & Fact-checker catches, revision loop \\
Model 404 (bad name) & All LLM calls fail & Fallback paths; fixed via config \\
Port 8000 busy & Django unreachable & Configurable BACKEND\_API\_URL \\
\bottomrule
\end{longtable}

\begin{warningbox}
\textbf{Resilience principle:} Har external call (LLM, search, vector store) try/except mein wrapped hai aur system \textbf{degraded mode} mein kaam karta rehta hai. Interview line: ``Maine fail-fast nahi, fail-graceful kiya --- user ko hamesha kuch na kuch output milta hai, error message ke saath.''
\end{warningbox}

\subsection{Scaling Story --- Kya Tootega, Kya Fix Hoga}

\begin{conceptbox}
\textbf{Scaling analysis (interview-ready):}

\begin{itemize}
\item \textbf{10 users:} Sab kuch chalta hai --- Streamlit single instance, SQLite, local ChromaDB.
\item \textbf{100 users:} SQLite concurrency limit $\rightarrow$ PostgreSQL. ChromaDB local $\rightarrow$ managed vector DB (Pinecone/Qdrant). Streamlit $\rightarrow$ multi-instance behind load balancer.
\item \textbf{10k documents:} Embedding batch jobs (offline), HNSW index tuning, collection sharding.
\item \textbf{1M requests:} Task queue (Celery/Redis) for async pipeline execution; results cached; webhook/poll for completion.
\item \textbf{Cost:} Groq cheap; embedding local free; Tavily paid --- caching + dedup search results.
\end{itemize}

\textbf{Core scalability insight:} Pipeline synchronous hai (176s/request). Production mein ise \textbf{async} karna first priority hai --- user request par 176s wait nahi karta.
\end{conceptbox}

\subsection{Why This Architecture (Not Monolith)}

\begin{interviewbox}
\textbf{Q: ``Ek hi Python app mein sab kuch kyun nahi? 3 interfaces kyun?''}

\textbf{Answer:} 3 alag \textbf{consumption patterns} hain: (1) \textbf{Human} real-time progress + download chahta hai $\rightarrow$ Streamlit. (2) \textbf{AI assistant} (Claude Desktop) tools ke through call karna chahta hai $\rightarrow$ MCP. (3) \textbf{Programmatic/benchmark} users REST API chahte hain $\rightarrow$ Django. Teeno ek hi compiled graph ko call karte hain --- logic ek jagah, \textbf{interfaces alag}. Ye ``single core, multiple adapters'' pattern hai --- jaise REST, GraphQL, gRPC sab ek service ke upar.
\end{interviewbox}
